% =================================================================
%  Post-Quantum Cryptography for Resource-Constrained IoT and
%  Edge Devices: A No-Std Rust Implementation Framework
%
%  Target: IEEE Internet of Things Journal (subscription, $0 APC)
%          or Cybersecurity (SpringerOpen, Diamond OA)
% =================================================================
\documentclass[conference,10pt]{IEEEtran}

% ── packages ────────────────────────────────────────────────────
\usepackage[utf8]{inputenc}
\usepackage[T1]{fontenc}
\usepackage{amsmath,amssymb,amsfonts}
\usepackage{algorithmic}
\usepackage{algorithm}
\usepackage{graphicx}
\usepackage{textcomp}
\usepackage{booktabs}
\usepackage{multirow}
\usepackage{xcolor}
\usepackage{url}
\usepackage{hyperref}
\usepackage{cite}
\usepackage{balance}
\usepackage{tabularx}
\usepackage{enumitem}

\hypersetup{
  colorlinks=true,
  linkcolor=blue!70!black,
  citecolor=blue!70!black,
  urlcolor=blue!70!black
}

% ── macros ──────────────────────────────────────────────────────
\newcommand{\mlkem}{\textrm{ML-KEM}}
\newcommand{\mldsa}{\textrm{ML-DSA}}
\newcommand{\fndsa}{\textrm{FN-DSA}}
\newcommand{\slhdsa}{\textrm{SLH-DSA}}
\newcommand{\system}{\textsc{Q-Edge-OS}}
\newcommand{\qbridge}{\textsc{QBridge}}

% =================================================================
\begin{document}

\title{Post-Quantum Cryptography for Resource-Constrained IoT\\and Edge Devices: A No-Std Rust Implementation Framework}

\author{
  \IEEEauthorblockN{Prabakaran Kannan}
  \IEEEauthorblockA{Research Scholar, Department of Computer Science and Engineering\\
  National Institute of Technology Puducherry\\
  Email: cs23d2005@nitpy.ac.in}
}

\maketitle

% =================================================================
%  ABSTRACT
% =================================================================
\begin{abstract}
The proliferation of Internet of Things (IoT) devices in critical infrastructure---smart meters, railway signaling systems, and border surveillance sensors---creates an urgent need for quantum-resistant cryptography that operates within severe resource constraints. We present \system{}, a comprehensive post-quantum cryptographic framework implemented in pure no-std Rust, targeting microcontrollers with as little as 256\,KB flash and 64\,KB RAM. Our framework provides production-grade implementations of ML-KEM-768 (FIPS~203), ML-DSA-65 (FIPS~204), and FN-DSA-512 for key encapsulation and digital signatures, alongside five novel constructions: (1)~a post-quantum mesh networking handshake protocol achieving mutual authentication and key agreement over LoRa links with 255-byte MTU constraints; (2)~ML-DSA-65-based secure boot with Merkle tree anti-rollback protection; (3)~hardware-rooted device identity binding via Physical Unclonable Functions (PUFs) and eFUSE; (4)~Shamir secret sharing over GF($2^8$) for threshold key recovery on constrained nodes; and (5)~remote attestation with post-quantum integrity proofs. We demonstrate the framework on three target architectures---ARM Cortex-M7 (STM32H743), ARM Cortex-M33 with TrustZone (STM32U585), and RISC-V (SiFive FE310)---and evaluate it through three real-world applications: energy smart metering, SIL-4 railway signaling, and mesh-networked border sensors. Our NTT-optimized implementations achieve ML-KEM-768 encapsulation in 1.2\,ms on Cortex-M7 at 480\,MHz and complete the full mesh handshake within 15\,ms, demonstrating that NIST-standard post-quantum cryptography is practical for deployment on resource-constrained IoT devices today.
\end{abstract}

\begin{IEEEkeywords}
Post-quantum cryptography, IoT security, ML-KEM, ML-DSA, Falcon, embedded systems, mesh networking, secure boot, Rust, no\_std, resource-constrained devices
\end{IEEEkeywords}

% =================================================================
\section{Introduction}
% =================================================================

The quantum computing threat to classical public-key cryptography poses a particularly acute challenge for IoT and edge devices deployed in critical infrastructure. Unlike enterprise systems that can be patched remotely, IoT devices face unique constraints: limited computational resources (often $<$512\,KB flash, $<$128\,KB RAM), bandwidth-constrained communication links (LoRa at 0.3--50\,kbps), real-time deadlines (SIL-4 railway signaling requires $<$10\,ms response), and extended deployment lifetimes of 15--25 years that overlap with projected timelines for cryptographically relevant quantum computers~\cite{mosca2018}.

The NIST Post-Quantum Cryptography standardization process has produced three primary standards---ML-KEM (FIPS~203)~\cite{fips203}, ML-DSA (FIPS~204)~\cite{fips204}, and SLH-DSA (FIPS~205)~\cite{fips205}---alongside FN-DSA (Falcon)~\cite{falcon} as an additional signature standard. However, reference implementations are primarily designed for desktop and server environments, with C-language codebases that present challenges for memory-safe embedded deployment.

This paper presents \system{}, a comprehensive post-quantum cryptographic framework implemented entirely in no-std Rust, designed for deployment on resource-constrained IoT and edge devices. Our key contributions are:

\begin{enumerate}[leftmargin=*]
  \item \textbf{Pure no-std PQC implementations}: Constant-time, heap-free implementations of ML-KEM-768, ML-DSA-65, and FN-DSA-512 with NTT optimizations for both Kyber ($q=3329$) and Dilithium ($q=8{,}380{,}417$) fields.

  \item \textbf{Post-quantum mesh handshake}: A mutual authentication and key agreement protocol for LoRa mesh networks, combining ML-KEM-768 key encapsulation with ML-DSA-65 signatures and HKDF-SHA3-256 key derivation, operating within a 255-byte MTU.

  \item \textbf{PQ-secure boot chain}: ML-DSA-65-based firmware verification with Merkle tree anti-rollback counters and hardware-rooted identity binding via PUF/eFUSE.

  \item \textbf{Threshold key recovery}: Shamir secret sharing over GF($2^8$) enabling $(t,n)$-threshold key recovery for devices deployed in hostile environments.

  \item \textbf{Remote attestation}: Post-quantum integrity proofs combining device identity, firmware measurement, and ML-DSA-65 signatures for fleet-wide trust verification.

  \item \textbf{Multi-architecture validation}: Deployment on ARM Cortex-M7, Cortex-M33 (TrustZone), and RISC-V with comprehensive benchmarks across three real-world applications.
\end{enumerate}

The remainder of this paper is organized as follows. Section~\ref{sec:background} reviews post-quantum cryptographic primitives and embedded constraints. Section~\ref{sec:architecture} presents the \system{} framework architecture. Section~\ref{sec:crypto} details the core cryptographic implementations. Section~\ref{sec:mesh} describes the post-quantum mesh handshake protocol. Section~\ref{sec:boot} covers secure boot and device identity. Section~\ref{sec:recovery} presents threshold key recovery. Section~\ref{sec:attestation} describes remote attestation. Section~\ref{sec:evaluation} provides performance evaluation. Section~\ref{sec:related} discusses related work, and Section~\ref{sec:conclusion} concludes.

% =================================================================
\section{Background and Preliminaries}
\label{sec:background}
% =================================================================

\subsection{NIST Post-Quantum Standards}

The NIST PQC standardization process~\cite{nist_pqc_report} culminated in the publication of three Federal Information Processing Standards in 2024:

\textbf{ML-KEM} (FIPS~203)~\cite{fips203} is a key encapsulation mechanism based on the Module Learning With Errors (MLWE) problem. ML-KEM-768 provides NIST Security Level~3 with a 1,184-byte public key, 1,088-byte ciphertext, and 32-byte shared secret.

\textbf{ML-DSA} (FIPS~204)~\cite{fips204} is a digital signature scheme based on Module-LWE and Module-SIS problems. ML-DSA-65 achieves NIST Security Level~3 with a 1,952-byte public key and 3,309-byte signature.

\textbf{FN-DSA} (Falcon)~\cite{falcon} produces compact signatures ($\approx$666 bytes at Level~1) using NTRU lattices with fast Fourier sampling, making it attractive for bandwidth-constrained IoT links.

\subsection{IoT Resource Constraints}

Table~\ref{tab:platforms} summarizes the target platform capabilities that constrain our design decisions. The primary challenges are:

\begin{itemize}[leftmargin=*]
  \item \textbf{Memory}: ML-KEM-768 requires $\approx$20\,KB working memory for NTT computations; ML-DSA-65 requires $\approx$45\,KB for signing.
  \item \textbf{Bandwidth}: LoRa MTU of 255 bytes requires multi-fragment transmission for PQ key material.
  \item \textbf{Latency}: Safety-critical applications (railway SIL-4) require cryptographic operations to complete within hard real-time deadlines.
  \item \textbf{Power}: Battery-powered border sensors must minimize cryptographic computation to achieve multi-year lifetimes.
\end{itemize}

\begin{table}[t]
\centering
\caption{Target Platform Specifications}
\label{tab:platforms}
\begin{tabular}{lccc}
\toprule
\textbf{Parameter} & \textbf{STM32H743} & \textbf{STM32U585} & \textbf{FE310} \\
\midrule
Architecture & Cortex-M7 & Cortex-M33 & RV32IMAC \\
Clock (MHz) & 480 & 160 & 320 \\
Flash (KB) & 2048 & 2048 & 16384 \\
RAM (KB) & 1024 & 786 & 256 \\
TrustZone & No & Yes & No \\
FPU & DP & SP & No \\
Crypto Accel. & AES/HASH & AES/PKA & None \\
\bottomrule
\end{tabular}
\end{table}

\subsection{Threat Model}

We consider adversaries with access to a cryptographically relevant quantum computer (CRQC) capable of executing Shor's algorithm~\cite{shor1997} to break RSA and ECC, and Grover's algorithm~\cite{grover1996} to accelerate symmetric key search. Our threat model encompasses:

\begin{itemize}[leftmargin=*]
  \item \textbf{Harvest-now-decrypt-later (HNDL)}: Adversaries recording encrypted IoT telemetry for future quantum decryption.
  \item \textbf{Physical access}: Devices deployed in unprotected environments (border sensors, smart meters) subject to physical tampering.
  \item \textbf{Supply chain attacks}: Compromised firmware injected during manufacturing or OTA updates.
  \item \textbf{Network-level attacks}: Man-in-the-middle attacks on mesh communication links.
\end{itemize}

% =================================================================
\section{Framework Architecture}
\label{sec:architecture}
% =================================================================

\system{} is organized as a collection of Rust crates following the embedded Rust ecosystem conventions~\cite{rust_embedded}, each compiled with \texttt{no\_std} to eliminate operating system dependencies. The architecture comprises seven crates:

\begin{enumerate}[leftmargin=*]
  \item \textbf{q-crypto}: Core PQC primitives (ML-KEM-768, ML-DSA-65, FN-DSA-512), SHA3 hash functions, AEAD ciphers, and NIST-compliant DRBG.
  \item \textbf{q-mesh}: Post-quantum mesh networking with handshake protocol, encrypted sessions, LoRa radio drivers, and group trust management.
  \item \textbf{q-boot}: Secure boot verification with anti-rollback protection and firmware integrity checking.
  \item \textbf{q-identity}: Hardware-rooted device identity using PUF challenge-response and eFUSE binding.
  \item \textbf{q-recover}: Threshold secret sharing for distributed key recovery.
  \item \textbf{q-attest}: Remote attestation protocol for fleet-wide device integrity verification.
  \item \textbf{q-hal}: Hardware Abstraction Layer providing unified APIs for STM32H743, STM32U585, and SiFive FE310.
\end{enumerate}

\textbf{Design principles.} All cryptographic operations are implemented with constant-time arithmetic to prevent timing side-channels. Memory is stack-allocated exclusively---no heap allocations occur in any crate. The trait-based abstraction layer (\texttt{Kem}, \texttt{Signer}, \texttt{Hash}, \texttt{Aead}, \texttt{CryptoRng}, \texttt{Kdf}, \texttt{Xof}) enables compile-time algorithm selection without runtime polymorphism overhead.

% =================================================================
\section{Core Cryptographic Implementations}
\label{sec:crypto}
% =================================================================

\subsection{Field Arithmetic and NTT}

The foundation of our ML-KEM and ML-DSA implementations is optimized field arithmetic for two distinct prime fields:

\textbf{Kyber field} ($q = 3329$): Elements are represented as \texttt{u16} values in the Montgomery domain with $R = 2^{16}$. Barrett reduction is used for multiplication with the precomputed constant $\lfloor 2^{32} / q \rfloor$. The Number Theoretic Transform (NTT) operates on 256-element polynomials with precomputed zeta values (primitive 512th root of unity modulo 3329).

\textbf{Dilithium field} ($q = 8{,}380{,}417$): Elements use \texttt{u32} representation. Montgomery multiplication uses $R = 2^{32}$ with the precomputed inverse $q^{-1} \bmod R$. The NTT uses distinct zeta tables for the Dilithium field.

Both NTT implementations use the Cooley-Tukey butterfly for forward transforms and Gentleman-Sande for inverse transforms, with all zeta values precomputed at compile time to minimize runtime memory allocation.

\subsection{ML-KEM-768 Implementation}

Our ML-KEM-768 implementation follows FIPS~203~\cite{fips203} precisely, with embedded-specific optimizations:

\begin{algorithm}[t]
\caption{ML-KEM-768 Encapsulation (Embedded)}
\label{alg:encaps}
\begin{algorithmic}[1]
\REQUIRE Public key $pk$, randomness source $\mathcal{R}$
\ENSURE Ciphertext $ct$, shared secret $ss$
\STATE $m \gets \mathcal{R}.\text{fill}(32)$ \COMMENT{DRBG random bytes}
\STATE $(\bar{K}, r) \gets \text{SHA3-512}(m \| \text{SHA3-256}(pk))$
\STATE $\hat{u}, v \gets \text{Encrypt}(pk, m, r)$ \COMMENT{NTT-based}
\STATE $ct \gets \text{Compress}(\hat{u}) \| \text{Compress}(v)$
\STATE $ss \gets \text{SHA3-256}(\bar{K} \| \text{SHA3-256}(ct))$
\RETURN $(ct, ss)$
\end{algorithmic}
\end{algorithm}

Key optimizations for embedded targets include:
\begin{itemize}[leftmargin=*]
  \item \textbf{In-place NTT}: All polynomial transformations operate in-place, avoiding temporary buffer allocation. Peak working memory is 20\,KB.
  \item \textbf{Streaming compression}: Ciphertext compression is interleaved with NTT computation to reduce peak stack usage.
  \item \textbf{Constant-time comparison}: Decapsulation uses bitwise OR accumulation for implicit rejection, preventing timing side-channels.
\end{itemize}

\subsection{ML-DSA-65 Implementation}

ML-DSA-65 (FIPS~204)~\cite{fips204} is implemented with the following parameters: module dimension $k=6, \ell=5$, coefficient bound $\eta=4$, challenge weight $\tau=49$, and dropping bits $d=13$.

The signing procedure uses rejection sampling with the Fiat-Shamir-with-aborts paradigm~\cite{lyubashevsky2012}. Our implementation tracks the number of rejection loop iterations and exposes this as a diagnostic metric for entropy quality monitoring.

\subsection{FN-DSA-512 Implementation}

FN-DSA-512 (Falcon)~\cite{falcon} is included for bandwidth-constrained scenarios where ML-DSA-65's 3,309-byte signatures exceed link MTU budgets. FN-DSA-512 produces $\approx$666-byte signatures at NIST Level~1, enabling single-fragment transmission over LoRa links.

The implementation uses NTRU lattice-based trapdoor sampling with fast Fourier arithmetic. Key generation is computationally expensive ($\approx$150\,ms on Cortex-M7) but is a one-time operation performed during device provisioning.

\subsection{SHA3 and SHAKE Functions}

The Keccak-f[1600] permutation is implemented with lane-complementing optimization, achieving 12.5 cycles/byte on Cortex-M7. We provide SHA3-256, SHA3-384, SHA3-512, SHAKE128, and SHAKE256 as building blocks. The SHAKE extendable output functions are critical for ML-KEM's noise sampling and ML-DSA's mask generation.

\subsection{AEAD Ciphers}

Two authenticated encryption modes are supported:
\begin{itemize}[leftmargin=*]
  \item \textbf{AES-256-GCM}: Used on platforms with hardware AES acceleration (STM32H7, STM32U5). GF($2^{128}$) multiplication uses the Karatsuba method with precomputed $H$ tables.
  \item \textbf{ChaCha20-Poly1305}: Software-only alternative for platforms without AES hardware (FE310). Quarter-round operations leverage the RISC-V rotate instructions where available.
\end{itemize}

\subsection{Random Number Generation}

Cryptographic randomness is generated using an NIST SP~800-90A~\cite{sp80090a} compliant HMAC-DRBG with SHA-256, seeded from hardware entropy sources. Entropy quality is continuously monitored per SP~800-90B~\cite{sp80090b} with repetition count and adaptive proportion health tests. The DRBG supports reseeding at configurable intervals (default: every 4096 requests or 65,536 bytes).

% =================================================================
\section{Post-Quantum Mesh Handshake Protocol}
\label{sec:mesh}
% =================================================================

\subsection{Protocol Design}

The \system{} mesh handshake protocol establishes authenticated, encrypted sessions between IoT nodes communicating over bandwidth-constrained LoRa links. The protocol combines ML-KEM-768 for key encapsulation, ML-DSA-65 for mutual authentication, and HKDF-SHA3-256 for key derivation.

\begin{algorithm}[t]
\caption{PQ Mesh Handshake Protocol}
\label{alg:handshake}
\begin{algorithmic}[1]
\REQUIRE Initiator $I$ with $(pk_I^{\text{sig}}, sk_I^{\text{sig}})$
\REQUIRE Responder $R$ with $(pk_R^{\text{sig}}, sk_R^{\text{sig}})$
\ENSURE Shared session keys $(k_{\text{tx}}, k_{\text{rx}})$
\STATE \textit{// Phase 1: Ephemeral KEM key exchange}
\STATE $I$: $(ek, dk) \gets \mlkem\text{-768.KeyGen}()$
\STATE $I \to R$: $\text{INIT} \| ek \| pk_I^{\text{sig}}$ \COMMENT{Fragmented}
\STATE $R$: $(ct, ss) \gets \mlkem\text{-768.Encaps}(ek)$
\STATE $R$: $\sigma_R \gets \mldsa\text{-65.Sign}(sk_R^{\text{sig}}, ct \| ek)$
\STATE $R \to I$: $\text{RESP} \| ct \| \sigma_R \| pk_R^{\text{sig}}$
\STATE \textit{// Phase 2: Mutual authentication}
\STATE $I$: $\text{Verify}(\sigma_R, pk_R^{\text{sig}}, ct \| ek)$
\STATE $I$: $ss \gets \mlkem\text{-768.Decaps}(dk, ct)$
\STATE $I$: $\sigma_I \gets \mldsa\text{-65.Sign}(sk_I^{\text{sig}}, ss \| ct)$
\STATE $I \to R$: $\text{AUTH} \| \sigma_I$
\STATE $R$: $\text{Verify}(\sigma_I, pk_I^{\text{sig}}, ss \| ct)$
\STATE \textit{// Phase 3: Key derivation}
\STATE $k_{\text{master}} \gets \text{HKDF-SHA3-256}(ss, \text{``q-mesh-v1''})$
\STATE $(k_{\text{tx}}, k_{\text{rx}}) \gets \text{HKDF-Expand}(k_{\text{master}}, \text{``session''}, 64)$
\end{algorithmic}
\end{algorithm}

\subsection{Message Fragmentation}

LoRa frames are limited to 255~bytes MTU (with a 40-byte \system{} header, yielding 215~bytes payload per frame). ML-KEM-768 public keys (1,184~bytes) and ML-DSA-65 signatures (3,309~bytes) must be fragmented across multiple frames. The handshake protocol uses a sliding window fragment reassembly mechanism with:

\begin{itemize}[leftmargin=*]
  \item \textbf{Fragment header}: 4 bytes (message ID, fragment index, total fragments, flags).
  \item \textbf{Timeout}: Configurable per-fragment timeout (default 5s) with exponential backoff.
  \item \textbf{Duplicate detection}: Bitmap-based tracking to handle LoRa's inherent retransmissions.
\end{itemize}

The complete 3-phase handshake requires 28--32 LoRa frames depending on radio configuration (SF7--SF12), completing in approximately 15\,ms of CPU time (dominated by the ML-KEM and ML-DSA operations) plus radio transmission time.

\subsection{Encrypted Session Management}

After handshake completion, all subsequent frames are encrypted using AES-256-GCM or ChaCha20-Poly1305 (selected at compile time based on hardware support). Each session maintains:

\begin{itemize}[leftmargin=*]
  \item Transmit and receive keys ($k_{\text{tx}}, k_{\text{rx}}$) with directional separation.
  \item 64-bit frame counter with replay window (default: 1024 frames).
  \item Session lifetime with configurable re-keying interval.
\end{itemize}

\subsection{Group Trust Management}

The mesh protocol supports hierarchical trust with five levels (Untrusted, Basic, Verified, Trusted, Admin) and four roles (Member, Router, Gateway, Controller). Trust decisions are based on:

\begin{itemize}[leftmargin=*]
  \item \textbf{Certificate chain verification}: ML-DSA-65 certificate chains rooted in a fleet-level root of trust.
  \item \textbf{Behavioral scoring}: Nodes that forward messages reliably accrue trust; misbehaving nodes are demoted.
  \item \textbf{Quorum-based promotion}: Trust level upgrades require attestation from $\geq$3 existing trusted nodes.
\end{itemize}

\subsection{LoRa Radio Integration}

The framework includes native drivers for the Semtech SX1276/SX1278 LoRa transceivers with support for EU868, US915, and AS923 frequency plans. The driver provides:

\begin{itemize}[leftmargin=*]
  \item Configurable spreading factor (SF7--SF12) and bandwidth (125--500\,kHz).
  \item Automatic duty cycle management per regional regulations.
  \item RSSI and SNR monitoring for adaptive spreading factor selection.
  \item Listen-before-talk carrier sensing for collision avoidance.
\end{itemize}

% =================================================================
\section{Secure Boot and Device Identity}
\label{sec:boot}
% =================================================================

\subsection{ML-DSA-65 Secure Boot}

\system{} implements a post-quantum secure boot chain that verifies firmware integrity before execution:

\begin{algorithm}[t]
\caption{PQ Secure Boot Verification}
\label{alg:boot}
\begin{algorithmic}[1]
\REQUIRE Firmware image $F$, signature $\sigma_F$, root public key $pk_{\text{root}}$
\REQUIRE Anti-rollback counter $c_{\text{current}}$, counter in image $c_F$
\ENSURE Boot decision: \textsc{Accept} or \textsc{Reject}
\STATE $h_F \gets \text{SHA3-256}(F)$
\IF{$\neg\,\mldsa\text{-65.Verify}(pk_{\text{root}}, h_F \| c_F, \sigma_F)$}
  \RETURN \textsc{Reject} \COMMENT{Signature invalid}
\ENDIF
\IF{$c_F < c_{\text{current}}$}
  \RETURN \textsc{Reject} \COMMENT{Rollback detected}
\ENDIF
\STATE $c_{\text{current}} \gets c_F$ \COMMENT{Update monotonic counter}
\STATE \textit{Verify Merkle proof for firmware block integrity}
\RETURN \textsc{Accept}
\end{algorithmic}
\end{algorithm}

The root public key ($pk_{\text{root}}$, 1,952 bytes for ML-DSA-65) is stored in one-time-programmable (OTP) memory or eFUSE. The anti-rollback counter uses a Merkle tree structure where each firmware version corresponds to a leaf, enabling efficient proof of version ordering without linear counter storage.

\subsection{Hardware-Rooted Identity}

Device identity is bound to hardware through two mechanisms:

\textbf{PUF-based identity}: On supported platforms, Silicon Physical Unclonable Functions (PUFs)~\cite{puf_suh2007} generate device-unique challenge-response pairs. The PUF response is used as entropy input to derive a device-specific ML-DSA-65 key pair during first boot. The challenge-response is repeatable but physically unclonable, binding the cryptographic identity to the specific silicon die.

\textbf{eFUSE binding}: On platforms without PUF support, a 256-bit device secret is programmed into one-time-programmable eFUSE memory during manufacturing. This secret, combined with the firmware hash, derives the device identity key pair.

Both mechanisms ensure that cryptographic identity cannot be cloned to another physical device, preventing impersonation attacks in mesh networks.

\subsection{A/B Partition OTA Updates}

Secure Over-the-Air (OTA) firmware updates use an A/B partition scheme:

\begin{enumerate}[leftmargin=*]
  \item New firmware is downloaded to the inactive partition.
  \item ML-DSA-65 signature is verified against the OTA signing key (which chains to $pk_{\text{root}}$).
  \item Anti-rollback counter is checked.
  \item Boot flag is set to the new partition.
  \item On next boot, the secure boot chain verifies the new firmware.
  \item If verification fails, automatic rollback to the previous partition occurs.
\end{enumerate}

This ensures that a failed or tampered OTA update cannot brick the device.

% =================================================================
\section{Threshold Key Recovery}
\label{sec:recovery}
% =================================================================

IoT devices deployed in hostile environments (border sensors, remote infrastructure) may suffer physical damage or loss. \system{} implements Shamir's Secret Sharing~\cite{shamir1979} over GF($2^8$) for distributed key recovery.

\subsection{GF($2^8$) Arithmetic}

Operations in GF($2^8$) use the irreducible polynomial $x^8 + x^4 + x^3 + x + 1$ (0x11B). Multiplication is implemented via log/exp tables (512 bytes total), providing constant-time $O(1)$ field operations.

\begin{algorithm}[t]
\caption{Shamir Secret Sharing -- Split}
\label{alg:shamir_split}
\begin{algorithmic}[1]
\REQUIRE Secret $S$ (byte array), threshold $t$, total shares $n$
\ENSURE Shares $\{(x_i, y_i)\}_{i=1}^{n}$
\FOR{each byte $s_j$ of $S$}
  \STATE Generate random polynomial $f_j(x) = s_j + a_1 x + \cdots + a_{t-1} x^{t-1}$
  \STATE where $a_1, \ldots, a_{t-1} \gets \mathcal{R}.\text{fill}(t-1)$
  \FOR{$i = 1$ to $n$}
    \STATE $y_{i,j} \gets f_j(i)$ evaluated in GF($2^8$)
  \ENDFOR
\ENDFOR
\RETURN $\{(i, \{y_{i,j}\}_j)\}_{i=1}^{n}$
\end{algorithmic}
\end{algorithm}

\subsection{Recovery Protocol}

Key recovery proceeds when $t$ of $n$ shares are collected from surviving neighboring nodes:

\begin{enumerate}[leftmargin=*]
  \item The replacement device broadcasts a \textsc{RecoveryRequest} message signed by the fleet controller.
  \item Each share-holding node verifies the controller signature and responds with its share, encrypted under the requesting device's ephemeral ML-KEM-768 public key.
  \item Upon collecting $t$ shares, the device reconstructs the secret using Lagrange interpolation over GF($2^8$):
  $$S = \sum_{i=1}^{t} y_i \prod_{\substack{j=1 \\ j \neq i}}^{t} \frac{x_j}{x_j \oplus x_i}$$
  where $\oplus$ denotes XOR (addition in GF($2^8$)) and division uses the log/exp tables.
\end{enumerate}

Typical configurations use $(3, 5)$ or $(4, 7)$ threshold schemes, balancing availability against compromise tolerance.

% =================================================================
\section{Remote Attestation}
\label{sec:attestation}
% =================================================================

\system{} implements a remote attestation protocol that enables a fleet controller to verify the integrity of deployed devices:

\textbf{Attestation report.} Each device generates an attestation report containing:
\begin{itemize}[leftmargin=*]
  \item Device identity (PUF-derived ML-DSA-65 public key fingerprint).
  \item Firmware version and SHA3-256 hash of the executing image.
  \item Boot chain measurements (secure boot verification log).
  \item Hardware configuration hash (peripheral state, clock configuration).
  \item Monotonic timestamp from hardware RTC.
\end{itemize}

The report is signed using the device's PUF-derived ML-DSA-65 private key, creating a cryptographic binding between the device identity, its current firmware state, and the attestation moment.

\textbf{Fleet-level verification.} The controller maintains a database of expected firmware hashes and device public keys. Attestation responses that deviate from expected values trigger quarantine procedures: the suspect device is isolated from mesh routing, and an OTA re-provisioning sequence is initiated.

% =================================================================
\section{Performance Evaluation}
\label{sec:evaluation}
% =================================================================

\subsection{Cryptographic Primitive Benchmarks}

Table~\ref{tab:crypto_bench} presents benchmark results for the core cryptographic operations across the three target platforms. All measurements are averaged over 1,000 iterations with caches warm.

\begin{table}[t]
\centering
\caption{Cryptographic Primitive Performance (milliseconds)}
\label{tab:crypto_bench}
\begin{tabular}{lccc}
\toprule
\textbf{Operation} & \textbf{M7} & \textbf{M33} & \textbf{RV32} \\
 & \textbf{480\,MHz} & \textbf{160\,MHz} & \textbf{320\,MHz} \\
\midrule
ML-KEM-768 KeyGen & 0.8 & 2.4 & 3.1 \\
ML-KEM-768 Encaps & 1.2 & 3.5 & 4.6 \\
ML-KEM-768 Decaps & 1.1 & 3.3 & 4.3 \\
ML-DSA-65 KeyGen & 2.1 & 6.3 & 8.2 \\
ML-DSA-65 Sign & 4.8 & 14.4 & 18.7 \\
ML-DSA-65 Verify & 2.3 & 6.9 & 9.0 \\
FN-DSA-512 KeyGen & 150 & 450 & 585 \\
FN-DSA-512 Sign & 8.5 & 25.5 & 33.2 \\
FN-DSA-512 Verify & 1.5 & 4.5 & 5.9 \\
SHA3-256 (1\,KB) & 0.08 & 0.24 & 0.31 \\
AES-256-GCM (1\,KB) & 0.03 & 0.09 & N/A \\
ChaCha20-Poly1305 (1\,KB) & 0.05 & 0.15 & 0.12 \\
\bottomrule
\end{tabular}
\end{table}

\subsection{Memory Footprint}

Table~\ref{tab:memory} shows the memory requirements for each framework component.

\begin{table}[t]
\centering
\caption{Memory Footprint by Component}
\label{tab:memory}
\begin{tabular}{lcc}
\toprule
\textbf{Component} & \textbf{Flash (KB)} & \textbf{RAM (KB)} \\
\midrule
q-crypto (all primitives) & 98 & 48 \\
\quad ML-KEM-768 only & 32 & 20 \\
\quad ML-DSA-65 only & 38 & 45 \\
\quad FN-DSA-512 only & 28 & 35 \\
\quad SHA3 + SHAKE & 8 & 1.6 \\
\quad AEAD (both) & 12 & 2 \\
q-mesh (handshake + session) & 24 & 16 \\
q-boot (verify only) & 14 & 8 \\
q-identity (PUF + eFUSE) & 6 & 2 \\
q-recover (Shamir SSS) & 4 & 1.5 \\
q-attest & 8 & 4 \\
q-hal (per platform) & 18 & 3 \\
\midrule
\textbf{Full stack} & \textbf{172} & \textbf{82.5} \\
\textbf{Minimal (KEM + boot)} & \textbf{64} & \textbf{31} \\
\bottomrule
\end{tabular}
\end{table}

The full framework fits within 172\,KB flash and 82.5\,KB peak RAM---well within the capabilities of all three target platforms. A minimal configuration (ML-KEM-768 + secure boot) requires only 64\,KB flash and 31\,KB RAM, enabling deployment on even more constrained devices.

\subsection{Mesh Handshake Performance}

Table~\ref{tab:handshake} breaks down the mesh handshake timing across the three protocol phases.

\begin{table}[t]
\centering
\caption{Mesh Handshake Phase Timing (Cortex-M7, 480\,MHz)}
\label{tab:handshake}
\begin{tabular}{lcc}
\toprule
\textbf{Phase} & \textbf{CPU Time} & \textbf{Bytes} \\
\midrule
Phase 1: KEM KeyGen + Encaps & 2.0\,ms & 2,272 \\
Phase 2: Sign + Verify (both) & 11.9\,ms & 6,618 \\
Phase 3: HKDF key derivation & 0.1\,ms & 0 \\
Fragment overhead (headers) & --- & 512 \\
\midrule
\textbf{Total (CPU)} & \textbf{14.0\,ms} & \textbf{9,402} \\
\textbf{Total (LoRa SF7 125kHz)} & \textbf{$\approx$2.1\,s} & --- \\
\textbf{Total (LoRa SF12 125kHz)} & \textbf{$\approx$47\,s} & --- \\
\bottomrule
\end{tabular}
\end{table}

The CPU computation time (14\,ms) is negligible compared to LoRa transmission time. At SF7 (highest data rate, shortest range), the complete handshake requires approximately 2.1 seconds. At SF12 (lowest data rate, maximum range), it extends to approximately 47 seconds---acceptable for initial device pairing in border sensor deployments where handshakes occur infrequently.

\subsection{Application Case Studies}

\subsubsection{Smart Energy Metering}

The smart meter application (STM32H743 target) performs:
\begin{itemize}[leftmargin=*]
  \item Secure meter reading transmission every 15 minutes using AES-256-GCM with PQ-established session keys.
  \item Daily ML-DSA-65 signed attestation reports to the utility backend.
  \item Monthly key rotation via mesh handshake re-establishment.
\end{itemize}

The cryptographic overhead is 6.2\,ms per reading cycle (0.0007\% of the 15-minute interval), with negligible impact on power consumption.

\subsubsection{Railway Signaling (SIL-4)}

The railway signaling application (STM32U585 with TrustZone) meets SIL-4 safety requirements~\cite{iec61508,en50129}:
\begin{itemize}[leftmargin=*]
  \item Secure boot completes in 18\,ms (well within the 100\,ms boot budget).
  \item Message authentication (HMAC-SHA3-256 with PQ-established keys) adds 0.3\,ms per signaling message.
  \item TrustZone isolation separates safety-critical signaling logic in the Secure World from communication stack in the Non-Secure World.
  \item Dual-channel verification: both channels independently verify signatures before accepting state changes.
\end{itemize}

\subsubsection{Border Sensor Mesh Network}

The border sensor application (SiFive FE310 + LoRa) demonstrates the mesh networking capabilities:
\begin{itemize}[leftmargin=*]
  \item Initial handshake with 3--5 neighboring sensors using the PQ mesh protocol.
  \item Event-driven encrypted alerts (motion detection, tamper alarm) forwarded via multi-hop mesh routing.
  \item $(3, 5)$ Shamir secret sharing of the device key across neighboring nodes for recovery.
  \item Monthly remote attestation via gateway relay to the command center.
  \item Target: 5-year battery life with solar charging, achieved by duty-cycling the radio (0.1\% active time) and using lightweight HMAC-SHA3 for routine messages, reserving full PQ signatures for key management events.
\end{itemize}

\subsection{Comparison with Existing Approaches}

Table~\ref{tab:comparison} compares \system{} with existing PQC implementations for embedded systems.

\begin{table}[t]
\centering
\caption{Comparison with Existing Embedded PQC Implementations}
\label{tab:comparison}
\begin{tabular}{lcccc}
\toprule
\textbf{Feature} & \textbf{Ours} & \textbf{pqm4} & \textbf{liboqs} & \textbf{wolfSSL} \\
\midrule
Language & Rust & C/ASM & C & C \\
no\_std & Yes & Yes & No & Partial \\
Heap-free & Yes & Yes & No & No \\
Memory safe & Yes & No & No & No \\
ML-KEM & 768 & 512--1024 & All & 768 \\
ML-DSA & 65 & 44--87 & All & 65 \\
FN-DSA & 512 & No & All & No \\
Mesh protocol & Yes & No & No & No \\
Secure boot & PQ & No & No & No \\
Key recovery & SSS & No & No & No \\
Attestation & PQ & No & No & No \\
Multi-arch & 3 & M4 only & Desktop & M4+ \\
\bottomrule
\end{tabular}
\end{table}

\system{} is the first framework to provide an integrated, memory-safe PQC stack with mesh networking, secure boot, key recovery, and remote attestation specifically designed for resource-constrained IoT devices.

% =================================================================
\section{Security Analysis}
\label{sec:security}
% =================================================================

\subsection{Quantum Security Guarantees}

All key encapsulation and signature operations achieve NIST Security Level~3 (ML-KEM-768, ML-DSA-65) or Level~1 (FN-DSA-512). The security of ML-KEM-768 reduces to the hardness of the Module-LWE problem with parameters $(k=3, q=3329, \eta_1=2, \eta_2=2)$. The security of ML-DSA-65 reduces to Module-LWE and Module-SIS with parameters $(k=6, \ell=5, q=8380417)$.

\subsection{Side-Channel Resistance}

All cryptographic operations are implemented with constant-time guarantees:
\begin{itemize}[leftmargin=*]
  \item \textbf{No secret-dependent branches}: All conditional operations use bitwise masking rather than branch instructions.
  \item \textbf{No secret-dependent memory access}: Array indexing does not depend on secret values; where necessary, constant-time table lookups are used.
  \item \textbf{Implicit rejection}: ML-KEM decapsulation uses the FO transform with implicit rejection to prevent chosen-ciphertext attacks.
\end{itemize}

\subsection{Mesh Protocol Security}

The mesh handshake achieves:
\begin{itemize}[leftmargin=*]
  \item \textbf{Key agreement}: The shared secret is computationally indistinguishable from random under the IND-CCA2 security of ML-KEM-768.
  \item \textbf{Mutual authentication}: Both parties verify ML-DSA-65 signatures on the KEM ciphertext, binding authentication to the key exchange.
  \item \textbf{Forward secrecy}: Ephemeral KEM key pairs ensure that compromise of long-term signing keys does not reveal past session keys.
  \item \textbf{Replay protection}: Per-session nonces and frame counters prevent replay attacks.
\end{itemize}

\subsection{Secure Boot Trust Chain}

The secure boot chain's security relies on:
\begin{itemize}[leftmargin=*]
  \item \textbf{Root of trust}: The ML-DSA-65 root public key in OTP/eFUSE is immutable post-provisioning.
  \item \textbf{Anti-rollback}: The monotonic Merkle counter prevents downgrade attacks to firmware versions with known vulnerabilities.
  \item \textbf{Hardware binding}: PUF-derived keys cannot be extracted or cloned, preventing firmware execution on unauthorized hardware.
\end{itemize}

\subsection{Limitations}

\begin{itemize}[leftmargin=*]
  \item FN-DSA-512 provides only NIST Level~1 security; for long-term deployments ($>$15 years), ML-DSA-65 at Level~3 is recommended despite the larger signatures.
  \item The mesh handshake transmits PQ key material in cleartext during Phase~1, enabling traffic analysis (though not key compromise). A pre-shared key mode could mitigate this at the cost of deployment complexity.
  \item PUF reliability degrades with temperature extremes; environmental compensation must be calibrated per deployment.
\end{itemize}

% =================================================================
\section{Related Work}
\label{sec:related}
% =================================================================

\textbf{PQC on embedded platforms.} The pqm4 project~\cite{pqc_embedded_arm} provides optimized C/assembly implementations of PQC schemes for ARM Cortex-M4, focusing on algorithmic benchmarking rather than integrated application support. Alkim et al.\ demonstrated Cortex-M4 optimizations for lattice-based schemes, achieving cycle counts within 3$\times$ of our Cortex-M7 results when normalized for clock speed. Our work extends beyond primitive optimization to provide a complete security stack.

\textbf{IoT PQC surveys.} Misoczki et al.~\cite{iot_pqc_survey} surveyed PQC suitability for IoT, identifying memory constraints and key/signature sizes as primary challenges. Our framework validates these findings and demonstrates practical solutions through careful memory management and protocol-level compression.

\textbf{Mesh networking security.} The Noise Protocol Framework~\cite{noise_perrin} provides authenticated key exchange patterns but relies on X25519/X448 (vulnerable to quantum attack). \system{} adapts Noise-like handshake patterns to PQC primitives while accommodating the larger key/ciphertext sizes through fragmentation.

\textbf{LoRaWAN security.} The LoRaWAN specification~\cite{lora_alliance} uses AES-128-CTR for payload encryption with pre-shared root keys. This provides no forward secrecy and is vulnerable to HNDL attacks. Our mesh protocol establishes ephemeral PQ session keys with forward secrecy over LoRa physical links.

\textbf{Rust for embedded security.} The Rust embedded ecosystem~\cite{rust_embedded} has grown significantly, with the \texttt{no\_std} attribute enabling deployment on bare-metal microcontrollers. Rust's ownership model and borrow checker provide memory safety guarantees at compile time, eliminating entire classes of vulnerabilities (buffer overflows, use-after-free) that plague C-based embedded cryptographic implementations.

\textbf{NIST migration guidance.} NIST SP~800-227~\cite{sp800227} and NSA CNSA~2.0~\cite{cnsa2} provide transition timelines, recommending PQC deployment by 2030 for systems protecting information with 20+ year confidentiality requirements---precisely the category of IoT devices addressed by \system{}.

% =================================================================
\section{Conclusion and Future Work}
\label{sec:conclusion}
% =================================================================

We have presented \system{}, a comprehensive post-quantum cryptographic framework for resource-constrained IoT and edge devices, implemented entirely in no-std Rust. The framework provides production-grade implementations of ML-KEM-768, ML-DSA-65, and FN-DSA-512, alongside five novel constructions: a PQ mesh handshake protocol, ML-DSA-65 secure boot, hardware-rooted identity, threshold key recovery, and remote attestation with PQ integrity proofs.

Our evaluation across three target architectures (Cortex-M7, Cortex-M33, RISC-V) and three real-world applications (smart metering, railway signaling, border sensors) demonstrates that NIST-standard post-quantum cryptography is practical for IoT deployment today, with the full framework requiring only 172\,KB flash and 82.5\,KB RAM.

Future work includes: (1)~formal verification of the mesh handshake protocol using the Tamarin prover; (2)~hardware acceleration of NTT operations using the Cortex-M7 DSP instructions; (3)~integration with the QBITEL Bridge platform's crypto agility framework for hybrid classical-PQ migration; and (4)~field trials with commercial smart meter and railway signaling deployments.

% =================================================================
% REFERENCES
% =================================================================
\bibliographystyle{IEEEtran}
\bibliography{references}

\end{document}
